% ============================================================================
%  Section 2 -- Calibration model
% ============================================================================
\section{Calibration Model}
\label{sec:method}

Sensor $i$ reports a raw reading $r_{i}$. We model the true temperature as a
linear function of that reading,
\begin{equation}
    T_{i} = a_{i}\, r_{i} + b_{i},
    \label{eq:model}
\end{equation}
where $a_{i}$ is the gain and $b_{i}$ is the offset. All sensors sit in a
stirred water bath next to a reference thermometer, and the bath steps through
set points between \qty{5}{\celsius} and \qty{45}{\celsius}.

\subsection{Least-squares fit}

With reference temperatures $T_{\mathrm{ref},k}$ at
\MarkEdit{edits}[math]{N = 2}{N \geq 3} set points, the parameters in
\cref{eq:model} follow from ordinary least squares~\cite{press2007numerical}:

\begin{MarkEnv}{changes}[math]
    (\hat{a}_{i}, \hat{b}_{i}) =
    \argmin_{a,\,b} \sum_{k=1}^{N}
    \bigl( T_{\mathrm{ref},k} - a\, r_{i,k} - b \bigr)^{2}
\end{MarkEnv}

\noindent
The root-mean-square residual of this fit estimates the remaining calibration
uncertainty of sensor~$i$. \Mark{changes}{A third set point makes this residual meaningful, because two points always fit a line exactly.}

\Cref{fig:fit} \Mark{missing}{TODO: describe the fitted lines once the figure exists.}

\MarkPlaceholderFigure{missing}{%
  Plot raw readings against the reference temperature for each sensor,
  with the fitted calibration lines.%
}{fig:fit}{Raw sensor readings against the reference temperature, with
per-sensor linear fits.}
